\documentclass[12pt,a4paper]{article}
\usepackage[utf8]{inputenc}
\usepackage[T1]{fontenc}
\usepackage{lmodern}
\usepackage[french]{babel}
\usepackage{amsmath,amssymb,amsthm}
\usepackage[top=2.5cm,bottom=2.5cm,left=2.8cm,right=2.8cm]{geometry}
\usepackage{parskip}
\usepackage{booktabs}
\usepackage{graphicx}
\usepackage{tikz}

\newtheorem{theorem}{Théorème}
\newtheorem{lemma}[theorem]{Lemme}
\newtheorem{corollary}[theorem]{Corollaire}
\newtheorem{proposition}[theorem]{Proposition}
\newtheorem{remark}[theorem]{Remarque}
\newtheorem{definition}[theorem]{Définition}
\newtheorem{example}[theorem]{Exemple}

\newcommand{\Inc}{\ensuremath{\mathrm{Inc}}}
\newcommand{\Exc}{\ensuremath{\mathrm{Exc}}}
\newcommand{\vbar}{\bar{v}}

\begin{document}

\begin{abstract}
Les machines de Tsetlin constituent une alternative interprétable et
efficace matériellement aux réseaux de neurones, mais leur analyse théorique
reste difficile en raison des interactions entre automates au sein des
clauses. Des approches récentes, telles que Probabilistic Concept Learning
(PCL), découplent les automates et introduisent, pour chaque clause, un
paramètre de transition $p_k\in(1/2,1)$ afin de rendre leur analyse
théorique possible. Contrairement à PCL, qui conserve des mécanismes de
mise à jour distincts pour les exemples positifs et négatifs, nous
introduisons une règle locale unique à quatre cas. Chaque automate associé
à un littéral est mis à jour uniquement à partir de la paire $(x_i,y)$,
sans dépendre de la sortie de la clause ni du vote de l'ensemble. Cette
structure permet de modéliser exactement chaque automate comme une chaîne
de Markov naissance--mort à $2N$ états.

Nous dérivons sa distribution stationnaire et caractérisons les régimes
d'inclusion et d'exclusion des littéraux sous une loi jointe arbitraire
$P(X,Y)$. Pour une clause comportant $n$ variables, dont $K$ sont fixées
par une condition, nous montrons que
\[
\limsup_{t\to\infty}
\mathbb{P}\!\left(C_t\neq C^\star\right)
\leq
2(n-K)\gamma_{\min}^{-N},
\]
où $C^\star$ désigne la configuration de littéraux stationnairement
préférée par la dynamique et où $\gamma_{\min}>1$ mesure le plus faible
biais d'inclusion ou d'exclusion parmi les automates. Nous intégrons ensuite
cette règle dans un classifieur one-vs-rest composé de clauses soumises à
des conditions aléatoires et agrégées par une pondération de type AdaBoost.
Nous donnons également des conditions explicites d'identifiabilité de la
conjonction génératrice sous dépendance et déséquilibre arbitraires des
features, et montrons que cette conjonction reste récupérable sous un bruit
de labels arbitraire de probabilité $\eta<\mu/(S+a)$. Sur le protocole de
récupération de conjonction de PCL (cible aléatoire, table de vérité
complète, sans bruit), notre règle récupère la cible exactement dans
$100/100$ essais jusqu'à $n{=}12$ littéraux, contre une dégradation nette
de PCL ($99\%$ à $n{=}4$, $37\%$ à $n{=}12$). Intégrée dans un classifieur
one-vs-rest conditionné et boosté par AdaBoost, la règle atteint par
ailleurs $100\%\pm0\%$ sur XOR bruité et des performances compétitives sur
Iris, Digits et MNIST (détail des budgets de calcul en Section~6). Les
garanties théoriques portent sur la règle locale sous une loi $P(X,Y)$
fixe ; l'ensemble conditionné et boosté évalué expérimentalement entraîne
chaque clause sous une distribution adaptative $Q_m$ qui change à chaque
round de boosting. Nous montrons que ces mêmes garanties s'appliquent telles
quelles à chaque round pris isolément (avec $P$ remplacée par $Q_m$), mais
la convergence du vote pondéré final combinant tous les rounds reste
ouverte.
\end{abstract}

\section{Introduction}

Les machines de Tsetlin (TM) apprennent des clauses conjonctives de
littéraux à l'aide d'un ensemble d'automates finis, un par littéral
candidat, mis à jour par un mécanisme de feedback de Type~I / Type~II.
Ce mécanisme couple les automates entre eux : la mise à jour de chaque
automate dépend non seulement de l'exemple courant, mais aussi de la
sortie de la clause tout entière et du vote global de l'ensemble. Ce
couplage rend l'analyse théorique de la dynamique de convergence
difficile : même pour une seule clause, les automates n'évoluent pas comme
des chaînes de Markov indépendantes, et la littérature manque encore de
garanties de convergence exactes sous une distribution de données
arbitraire.

Probabilistic Concept Learning (PCL) est l'approche la plus proche de la
nôtre : elle découple les automates et associe à chaque clause un unique
paramètre de transition $p_k\in(1/2,1)$, ce qui permet d'obtenir un résultat
de convergence qualitatif. Cette simplification reste néanmoins partielle :
PCL conserve des règles de mise à jour distinctes selon que l'exemple est
positif ou négatif, et son analyse ne caractérise pas exactement la loi
stationnaire de chaque automate sous une distribution de données
arbitraire.

Nous proposons une règle locale différente, à quatre cas, où la mise à
jour d'un automate ne dépend que de la paire $(x_i,y)$ -- ni de la sortie
de la clause, ni des autres automates, ni du vote de l'ensemble. Ce
découplage total permet de modéliser exactement chaque automate comme une
chaîne de Markov naissance--mort réfléchie à $2N$ états, et d'en dériver
la loi stationnaire exacte sous une loi jointe $P(X,Y)$ arbitraire.

\textbf{Contributions.} (1) Une règle de mise à jour locale à quatre cas,
découplée de la sortie de la clause et du vote de l'ensemble. (2) La loi
stationnaire exacte de chaque automate sous une loi $P(X,Y)$ arbitraire, et
la caractérisation des trois régimes de sélection d'un littéral qui en
découlent. (3) Une borne de concentration structurelle multi-features en
$2(n-K)\gamma_{\min}^{-N}$ pour une clause entière. (4) Des conditions
explicites d'identification de la conjonction génératrice sous dépendance
et déséquilibre arbitraires des features, et sa récupération sous un bruit
de labels arbitraire de probabilité $\eta<\mu/(S+a)$. (5) Un classifieur
one-vs-rest pratique combinant cette règle, des conditions aléatoires et un
poids de type AdaBoost, évalué sur XOR bruité, Iris, Digits et MNIST.

\section{Contexte et travaux connexes}

\subsection{Machines de Tsetlin}

Une clause TM est une conjonction de littéraux $x_i$ ou $\neg x_i$,
chacun gouverné par un automate à $2N$ états dont la moitié inférieure
correspond à l'action Exclude et la moitié supérieure à l'action Include.
Le feedback dépend du label $y$, de la sortie de la clause, et -- pour le
Type~II en particulier -- du vote global de l'ensemble sur cet exemple :
le Type~I renforce les littéraux cohérents avec un exemple positif que la
clause reconnaît, le Type~II pénalise les littéraux inclus à tort sur un
exemple négatif. Notre seul objectif ici est de préciser ce que notre règle
retire à ce mécanisme : la dépendance à la sortie de la clause et au vote de
l'ensemble.

\subsection{Probabilistic Concept Learning}

PCL découple les automates d'une même clause et leur associe, pour chacune,
un paramètre de transition unique $p_k\in(1/2,1)$ : la mise à jour de
chaque automate suit une marche aléatoire biaisée par $p_k$, avec des
règles différentes selon que l'exemple observé est positif ($y=1$) ou
négatif ($y=0$) (mécanisme vérifié directement dans le code officiel des
auteurs, \texttt{dpcl-classifier/Main\_code\_AAAI25}). Sous cette hypothèse,
PCL établit un résultat de convergence garantissant la récupération d'une
conjonction cible.

PCL conserve une structure de feedback héritée des machines de Tsetlin et
fournit un résultat de convergence qualitatif sous des hypothèses
spécifiques sur $p_k$. Nous cherchons ici une caractérisation exacte, à
états finis, de la dynamique de chaque automate sous une distribution de
données $P(X,Y)$ arbitraire -- sans restreindre a priori la forme de cette
distribution.

\begin{center}
\begin{tabular}{lccc}
\toprule
Propriété & TM & PCL & Notre règle \\
\midrule
Dépend de la sortie de clause & Oui & Non & Non \\
Dépend du vote global & Oui & Non & Non \\
Paramètre par clause & $(s,T)$ & $p_k$ & $(S,a)$ \\
Loi stationnaire finie explicite & Non & Non & Oui \\
Loi $P(X,Y)$ arbitraire & Non démontré & Non & Oui \\
\bottomrule
\end{tabular}
\end{center}

\section{Règle d'apprentissage locale à quatre cas}

\subsection{Cadre et notations}

Soit $(X_t,Y_t)_{t\ge1}\overset{\mathrm{i.i.d.}}{\sim}P$ de loi jointe
$P(X,Y)$ arbitraire, avec $X=(X_1,\ldots,X_n)\in\{0,1\}^n$ et $Y\in\{0,1\}$.
Pour chaque variable $X_i$, deux automates à $2N$ états gouvernent le
littéral $x_i$ et sa négation $\neg x_i$ : $v_i$ et $\vbar_i$, avec
\[
  \mathrm{Include}(s) = \mathbf 1[s>N], \qquad s\in\{1,\ldots,2N\}.
\]
La clause induite par ces automates est
\[
  C(x) = \bigwedge_{i:\,v_i>N} x_i \;\wedge\; \bigwedge_{i:\,\vbar_i>N} \neg x_i.
\]
Les mouvements élémentaires d'un automate sont
\[
  s \xrightarrow{\ \Inc\ } \min(s+1,2N), \qquad
  s \xrightarrow{\ \Exc\ } \max(s-1,1)
\]
(bornes réfléchissantes : un automate déjà dans l'état $2N$ (resp.\ $1$) reste
sur place s'il reçoit à nouveau l'ordre $\Inc$ (resp.\ $\Exc$) -- d'où la
chaîne de Markov naissance--mort \emph{réfléchie} annoncée en Section~3.1).

\subsection{Mise à jour à quatre cas}

\begin{center}
\begin{tabular}{cccc}
\toprule
$x_i$ & $y$ & Action sur $v_i$ & Action sur $\vbar_i$ \\
\midrule
0 & 0 & vers $\Inc$ avec proba.\ $S$ & vers $\Exc$ avec proba.\ $S$ \\
0 & 1 & vers $\Exc$ avec proba.\ $a$ & pas de mise à jour \\
1 & 0 & vers $\Exc$ avec proba.\ $S$ & vers $\Inc$ avec proba.\ $S$ \\
1 & 1 & pas de mise à jour & vers $\Exc$ avec proba.\ $a$ \\
\bottomrule
\end{tabular}
\end{center}
avec $S,a\in(0,1]$. Contrairement aux mécanismes de feedback des TM et de
PCL, cette mise à jour ne dépend que de la paire locale $(x_i,y)$, et non
de la sortie de la clause, du vote de l'ensemble, ou de l'état des autres
automates.

\subsection{Interprétation}

Sur un exemple négatif ($y=0$), la règle pousse chaque littéral vers l'état
qui le rend \emph{violé} par cet exemple (i.e.\ qui empêche la clause de
répondre $1$ à tort) : $S$ contrôle la force de cet apprentissage. Sur un
exemple positif ($y=1$), la règle exclut les littéraux \emph{incompatibles}
avec cet exemple (ceux que l'exemple positif viole) : $a$ contrôle la force
de cette exclusion. À ce stade, rien ne garantit encore que la structure
ainsi sélectionnée coïncide avec une éventuelle conjonction génératrice
sous-jacente -- ce point est traité séparément en Section~4.

\subsection{Représentation graphique}

La Table~1 (Section~3.2) reste la représentation formelle de référence. La
Figure~\ref{fig:automate} n'en donne qu'une lecture graphique : la chaîne
d'états d'un automate, avec la frontière Exclude/Include entre $N$ et $N{+}1$.

\begin{figure}[h]
\centering
\begin{tikzpicture}[>=stealth, every node/.style={font=\small}]
  \node[draw, circle, minimum size=7mm] (s1)  at (0,0)   {$1$};
  \node[draw, circle, minimum size=7mm] (s2)  at (1.6,0) {$2$};
  \node[draw, circle, minimum size=7mm] (sN)  at (4.4,0) {$N$};
  \node[draw, circle, minimum size=7mm] (sN1) at (6.4,0) {$N{+}1$};
  \node[draw, circle, minimum size=7mm] (s2N) at (9.2,0) {$2N$};
  \node at (3.0,0) {$\cdots$};
  \node at (7.8,0) {$\cdots$};
  \draw[->] (s1) -- (s2);
  \draw[->] (s2) -- (3.4,0);
  \draw[->] (3.6,0) -- (sN);
  \draw[->] (sN) -- (sN1);
  \draw[->] (sN1) -- (7.2,0);
  \draw[->] (7.4,0) -- (s2N);
  \draw[dashed] (5.4,0.9) -- (5.4,-0.9);
  \node at (2.2,-1.1) {\footnotesize Exclude};
  \node at (7.8,-1.1) {\footnotesize Include};
\end{tikzpicture}
\caption{Chaîne d'états d'un automate à $2N$ états. Les flèches représentent
les mouvements $\Inc$ et $\Exc$ (Section~3.1) ; la frontière (pointillés)
sépare Exclude ($1,\ldots,N$) et Include ($N{+}1,\ldots,2N$).}
\label{fig:automate}
\end{figure}

\subsection{Algorithme d'apprentissage local}

\begin{quote}\ttfamily\small
\textbf{Entrée} : échantillons $(x_t,y_t)$, $2N$ états, paramètres $S,a$.\\[4pt]
initialiser $(v_i,\vbar_i)$ pour chaque feature $i$\\
\textbf{pour} $t = 1$ à $T$ :\\
\hspace*{1.5em}\textbf{pour} $i = 1$ à $n$ :\\
\hspace*{3em}mettre à jour $(v_i,\vbar_i)$ selon la Table~1, à partir de $(x_{t,i},y_t)$\\
\textbf{retourner} la clause induite par les littéraux inclus
\end{quote}

À ce stade, ni le boosting AdaBoost, ni les conditions aléatoires, ni le
classifieur multiclasse ne sont encore introduits : ils font l'objet de la
Section~5.

\section{Analyse théorique}

Cette section prouve la convergence de la règle locale à 4 cas elle-même,
et non celle du classifieur complet ; son intégration dans un ensemble
conditionné et boosté est traitée séparément en Section~5, comme une
application dont l'analyse de convergence complète reste un travail futur.

\subsection{Dynamique d'un automate isolé}

Pour $v_i$, les probabilités de transition intérieures s'écrivent, pour
toute loi $P$ :
\begin{align}
  p_i^+ &= S\,P(X_i{=}0,Y{=}0), & p_i^- &= S\,P(X_i{=}1,Y{=}0) + a\,P(X_i{=}0,Y{=}1); \label{eq:pv}\\
  \bar p_i^+ &= S\,P(X_i{=}1,Y{=}0), & \bar p_i^- &= S\,P(X_i{=}0,Y{=}0) + a\,P(X_i{=}1,Y{=}1). \label{eq:pvbar}
\end{align}
On note $\rho_i=p_i^+/p_i^-$, $\bar\rho_i=\bar p_i^+/\bar p_i^-$. Pour
$p_i^+,p_i^->0$, chaque marginal $(v_{i,t})_t$ est une chaîne de Markov
naissance--mort réfléchie sur $\{1,\ldots,2N\}$.

\subsection{Loi stationnaire exacte}

\begin{theorem}\label{thm:stat}
Soit $\rho=p^+/p^-$, avec $p^+,p^->0$ (indices $i$ omis). La chaîne $(v_t)$
est irréductible et apériodique sur $\{1,\ldots,2N\}$, donc
$\mathcal L(v_t)\xrightarrow[t\to\infty]{}\pi$ (convergence en loi vers
l'unique loi stationnaire, théorème ergodique standard pour les chaînes de
Markov finies irréductibles apériodiques), de loi stationnaire
\[
  \pi(k) = \frac{(1-\rho)\rho^{k-1}}{1-\rho^{2N}}\ (\rho\ne1), \qquad
  \pi(k) = \frac1{2N}\ (\rho=1), \qquad
  \pi(\Inc) = \frac{\rho^N}{1+\rho^N}.
\]
Si $\rho>1$ : $\pi(\Inc)\to1$ et $1-\pi(\Inc)\le\rho^{-N}$. Si $\rho<1$ :
$\pi(\Exc)\to1$ symétriquement. De plus $\frac1T\sum_{t\le T}\mathbf1[v_t>N]\to\pi(\Inc)$ p.s.
\emph{Attention} : ceci est une convergence en loi (stationnaire) et une
convergence p.s.\ de la \emph{fréquence temporelle} -- pas une convergence
p.s.\ de $v_t$ vers un état fixe. Dans le cas ergodique ($p^+,p^->0$),
l'automate continue de se déplacer indéfiniment ; il n'y a stabilisation
permanente que dans le cas absorbant (Corollaire~\ref{cor:absorb}).
\end{theorem}

\begin{proof}
\emph{Irréductibilité.} Chaîne de naissance-mort à transitions intérieures
strictement positives.

\emph{Noyau explicite.} Pour $2\le k\le 2N-1$ : $P(k,k{+}1)=p^+$,
$P(k,k{-}1)=p^-$, $P(k,k)=1-p^+-p^-$. Aux bornes réfléchissantes :
$P(1,1)=1-p^+$, $P(1,2)=p^+$, $P(2N,2N)=1-p^-$, $P(2N,2N{-}1)=p^-$.

\emph{Apériodicité.} $P(1,1)=1-p^+\ge p^->0$ : l'état $1$ a un cycle de
longueur $1$, donc une période $d(1)=1$. Par irréductibilité, la période est
une propriété de classe : toute la chaîne est apériodique.

\emph{Forme géométrique.} Bilan détaillé $\pi(k)p^+=\pi(k{+}1)p^-$ donne
$\pi(k{+}1)/\pi(k)=\rho$ constant, soit $\pi(k)=\pi(1)\rho^{k-1}$ ;
normalisation $\sum_{k=1}^{2N}\pi(k)=1=\pi(1)\sum_{k=0}^{2N-1}\rho^k$ donne
$\pi(1)=(1-\rho)/(1-\rho^{2N})$ (série géométrique finie), d'où les formules
annoncées. $\pi(\Inc)=\sum_{k=N+1}^{2N}\pi(k)=\pi(1)\rho^N\sum_{k=0}^{N-1}\rho^k
=\pi(1)\rho^N(1-\rho^N)/(1-\rho)=\rho^N(1-\rho^N)/(1-\rho^{2N})=\rho^N/(1+\rho^N)$.
La convergence p.s.\ de la moyenne temporelle est le théorème ergodique
standard.
\end{proof}

\begin{corollary}[Cas absorbant]\label{cor:absorb}
Si $p^-=0,p^+>0$ : absorption p.s.\ en $2N$ en temps fini,
$\mathbb E[T_{2N}]=(2N-s_0)/p^+$ ; première entrée dans $\Inc$ (franchissement
de $N$), généralement plus rapide, $\mathbb E[T_{\Inc}]=(N+1-s_0)_+/p^+$. Si
$p^+=0,p^->0$ : symétrique en $1$/$\Exc$. Si $p^+=p^-=0$ : aucune transition
ne se produit jamais (chaque état trivialement absorbant), l'automate reste
figé à $s_0$.
\end{corollary}

\subsection{Régimes de sélection des littéraux}

\begin{lemma}\label{lem:excl}
Soit $\Delta=p^+-p^-$, $\bar\Delta=\bar p^+-\bar p^-$. Pour toute loi $P(X,Y)$ :
\[
  \Delta+\bar\Delta = -a\,P(Y{=}1) \le 0.
\]
Si $P(Y{=}1)>0$, l'inégalité est \emph{stricte} : $\Delta+\bar\Delta<0$, donc
$\Delta=\bar\Delta=0$ simultanément est impossible. Si $P(Y{=}1)=0$ (cas
dégénéré), $\Delta+\bar\Delta=0$ exactement.
\end{lemma}

\begin{proof}
De \eqref{eq:pv}--\eqref{eq:pvbar} : $\Delta = S\,P(0,0)-S\,P(1,0)-a\,P(0,1)$,
$\bar\Delta = S\,P(1,0)-S\,P(0,0)-a\,P(1,1)$. La somme élimine les termes en
$S$ : $\Delta+\bar\Delta=-a[P(0,1)+P(1,1)]=-a\,P(Y{=}1)\le0$, avec égalité
ssi $P(Y{=}1)=0$ (puisque $a>0$).
\end{proof}

\begin{corollary}\label{cor:asym}
Si $\Delta>0$ alors $\bar\Delta=-a\,P(Y{=}1)-\Delta<0$ strictement.
\end{corollary}

\begin{remark}
Le lemme exclut $\Delta>0$ et $\bar\Delta>0$ simultanément ; il n'exclut pas
$\Delta<0$ et $\bar\Delta<0$ simultanément (double exclusion, régime neutre).
Pour $N$ fini et $\rho\ne1$, l'état instantané peut néanmoins se trouver du
mauvais côté avec probabilité $\le\rho^{-N}\to0$.
\end{remark}

\begin{theorem}\label{thm:sel}
Supposons $\Delta\ne0$ et $\bar\Delta\ne0$. Notons $\gamma=\max(\rho,1/\rho)>1$,
$\bar\gamma=\max(\bar\rho,1/\bar\rho)>1$ (probabilité du mauvais côté toujours
$\le\gamma^{-N}$, resp.\ $\bar\gamma^{-N}$, quel que soit le signe de
$\Delta,\bar\Delta$). Alors exactement un des trois régimes suivants a lieu :
\begin{itemize}
  \item $\Delta>0$ (donc $\bar\Delta<0$) $\iff S\,P(0,0)>S\,P(1,0)+a\,P(0,1)$ :
  \[
    \lim_{N\to\infty}\;\limsup_{t\to\infty}\;\mathbb P\bigl((v_t,\vbar_t)\ne(\Inc,\Exc)\bigr) = 0.
  \]
  \item $\bar\Delta>0$ (donc $\Delta<0$) $\iff S\,P(1,0)>S\,P(0,0)+a\,P(1,1)$ :
  \[
    \lim_{N\to\infty}\;\limsup_{t\to\infty}\;\mathbb P\bigl((v_t,\vbar_t)\ne(\Exc,\Inc)\bigr) = 0.
  \]
  \item $\Delta<0$ et $\bar\Delta<0$ :
  \[
    \lim_{N\to\infty}\;\limsup_{t\to\infty}\;\mathbb P\bigl((v_t,\vbar_t)\ne(\Exc,\Exc)\bigr) = 0.
  \]
\end{itemize}
\end{theorem}

\begin{proof}
Les équivalences découlent de \eqref{eq:pv}--\eqref{eq:pvbar}. L'exclusion
mutuelle des trois régimes découle du Corollaire~\ref{cor:asym}. Passage du
marginal au joint (cas $\Delta>0$ : $\rho>1$ donc $\gamma=\rho$, $\bar\rho<1$
donc $\bar\gamma=1/\bar\rho$ ; similaire ailleurs), \emph{sans} invoquer de loi
stationnaire jointe (la chaîne $(v_t,\vbar_t)$ n'est pas nécessairement
irréductible -- ex.\ $S{=}1$, $P(Y{=}1){=}0$ conserve $v_t+\vbar_t$). En
prenant d'abord $\limsup_{t\to\infty}$ (la borne du Théorème~\ref{thm:stat}
porte sur la loi \emph{stationnaire}, atteinte seulement asymptotiquement en
$t$ -- fausse pour un $t$ petit et une initialisation arbitraire), puis par
union bound,
\[
  \limsup_{t\to\infty}\mathbb P\bigl((v_t,\vbar_t)\ne(\Inc,\Exc)\bigr)
  \le \limsup_{t\to\infty}\mathbb P(v_t\ne\Inc)+\limsup_{t\to\infty}\mathbb P(\vbar_t\ne\Exc)
  \le \gamma^{-N}+\bar\gamma^{-N},
\]
sans hypothèse d'indépendance entre $v_t$ et $\vbar_t$. Puis
$\gamma^{-N}+\bar\gamma^{-N}\xrightarrow[N\to\infty]{}0$ donne la double
limite annoncée. Noter $\bar\gamma^{-N}=\bar\rho^N$ ici (pas $\bar\rho^{-N}$,
puisque $\bar\rho<1$) : c'est pour éviter cette confusion de signe que
$\gamma=\max(\rho,1/\rho)$ est défini une fois pour toutes.
\end{proof}

\begin{figure}[h]
\centering
\includegraphics[width=0.7\textwidth]{figures/graphe_N_thm6.png}
\caption{Validation empirique du Théorème~\ref{thm:sel} : pour une paire
$(v,\bar v)$ sous une loi $P(X,Y)$ fixée ($\Delta>0$), probabilité
$P((v_t,\bar v_t)\ne(\Inc,\Exc))$ en fonction de $N$, comparée à la borne
exacte $\gamma^{-N}+\bar\gamma^{-N}$. Simulation fidèle à la règle réelle :
un seul exemple $(x,y)$ tiré par pas de temps (partagé par $v$ et $\bar v$,
comme dans le code), chaque automate concerné appliquant ensuite son propre
tirage indépendant -- $v$ et $\bar v$ ne sont donc pas simulés comme des
marches indépendantes. Nombre de répétitions adapté à chaque $N$ pour
garder au moins $\sim\!20$ évènements « mauvais régime » observés ; la
courbe empirique reste sous la borne sur toute la plage $N{=}2$ à $20$
(à $N{=}24$, $0$ évènement observé sur $25\,000$ répétitions, cohérent avec
une borne de l'ordre de $8\times10^{-9}$).}
\label{fig:thm6}
\end{figure}

\subsection{Concentration structurelle multi-features}

\begin{theorem}\label{thm:multi}
Clause à $n$ features, $K$ fixées par la condition, $d=n-K$ libres gouvernées par
$2d$ automates. Pour chaque automate $j$ (préférence stationnaire non
dégénérée), $r_j=p_j^+/p_j^-$, $\gamma_j=\max(r_j,1/r_j)>1$,
$\gamma_{\min}=\min_j\gamma_j$. Alors, sans hypothèse d'indépendance,
l'erreur de \emph{récupération structurelle} vérifie :
\[
  \limsup_{t\to\infty}\mathbb P(C_t\ne C^\star) \le \sum_{j=1}^{2d}\frac1{1+\gamma_j^N}
  \le 2(n-K)\,\gamma_{\min}^{-N} \xrightarrow[N\to\infty]{} 0.
\]
\end{theorem}

\begin{proof}
Soit $A_j$ l'évènement « l'automate $j$ est du mauvais côté ». Si tous les
$A_j$ sont faux, chaque automate est dans son état préféré, donc $C_t=C^\star$ ;
par contraposée, $\{C_t\ne C^\star\}\subseteq A_1\cup\cdots\cup A_{2d}$, d'où
par l'union bound (Boole) :
\[
  \mathbb P(C_t\ne C^\star)\le\sum_{j=1}^{2d}\mathbb P(A_j).
\]
Par le Théorème~\ref{thm:stat}, $\mathbb P(A_j)=\pi(\Exc)=1/(1+r_j^N)$ si
$r_j>1$, ou $\pi(\Inc)=r_j^N/(1+r_j^N)$ si $r_j<1$ -- dans les deux cas
exactement $1/(1+\gamma_j^N)$, d'où la première inégalité. Ensuite,
$1/(1+\gamma_j^N)<\gamma_j^{-N}$ (dénominateur plus grand), et
$\gamma_j\ge\gamma_{\min}$ donne $\gamma_j^{-N}\le\gamma_{\min}^{-N}$ ; en
remplaçant chacun des $2d$ termes par cette même borne on obtient
$2(n-K)\,\gamma_{\min}^{-N}$.
\end{proof}

\begin{remark}
Cas $r_j=1$ et $p_j^+=p_j^-=0$ exclus (non-dégénérescence) ; automate
absorbant (Corollaire~\ref{cor:absorb}) couvert par $\gamma_j=\infty$. La
borne utilise le pire automate, pas le meilleur. $C_t\ne C^\star$ est un
\emph{mismatch structurel} (deux structures différentes peuvent calculer la
même fonction logique par coïncidence), pas nécessairement une erreur de
classification.
\end{remark}

\begin{figure}[h]
\centering
\includegraphics[width=0.7\textwidth]{figures/graphe_N_thm7.png}
\caption{Validation empirique du Théorème~\ref{thm:multi} : erreur de
récupération structurelle $P(C_t\ne C^\star)$ en fonction de $N$, pour une
clause à 3 features de difficultés différentes ($\gamma_{\min}\approx1.17$).
Simulation fidèle à la règle réelle : un seul exemple $(x_1,x_2,x_3,y)$
tiré par pas de temps (loi jointe à indépendance conditionnelle en $Y$),
chacun des $6$ automates étant mis à jour à partir de ce même tirage
partagé -- pas $2d$ marches indépendantes. L'empirique reste toujours sous
la borne théorique sur toute la plage $N{=}2$ à $40$.}
\label{fig:thm7}
\end{figure}

\subsection{Récupération d'une conjonction cible}

Le Théorème~\ref{thm:sel} caractérise la structure préférée sous une loi $P$
arbitraire, mais ne dit pas si cette structure est la véritable cible. On le
montre ici pour une famille explicite de concepts : une conjonction cible,
sous une loi jointe $(X_1,\ldots,X_n,Y)$ \emph{quelconque} (features
corrélées, déséquilibrées, dépendantes) telle que $Y=Y^\star$.

\begin{definition}\label{def:ctrue}
Soit $C_{\mathrm{true}}(x)=\bigwedge_{j\in J_+}x_j\wedge\bigwedge_{j\in J_-}(1-x_j)$
une conjonction de $m=|J_+|+|J_-|\ge1$ littéraux, $J_+\cap J_-=\varnothing$
-- une fonction booléenne fixe, déterministe. On pose $Y^\star=C_{\mathrm{true}}(X)$,
la variable aléatoire obtenue en évaluant cette fonction sur $X$.
\end{definition}

\begin{theorem}[Identification sous dépendance de features arbitraire]\label{thm:identif}
Avec $C_{\mathrm{true}},Y^\star$ comme dans la Définition~\ref{def:ctrue},
supposons seulement $Y=Y^\star$ (p.s.) -- la loi jointe de
$(X_1,\ldots,X_n)$ est quelconque (dépendante, déséquilibrée). Pour chaque
variable $j$, notons
\[
  A_j=P(X_j{=}0,Y{=}0),\quad B_j=P(X_j{=}1,Y{=}0),\quad
  C_j=P(X_j{=}0,Y{=}1),\quad D_j=P(X_j{=}1,Y{=}1).
\]
Alors :
\begin{itemize}
  \item pour $j\in J_+$ : si $A_j>B_j$, alors $\Delta_j>0$ et $\bar\Delta_j<0$ ;
  \item pour $j\in J_-$ : si $B_j>A_j$, alors $\bar\Delta_j>0$ et $\Delta_j<0$ ;
  \item pour $k\notin J_+\cup J_-$ : si $-\tfrac aS D_k < A_k-B_k < \tfrac aS C_k$,
  alors $\Delta_k<0$ et $\bar\Delta_k<0$.
\end{itemize}
Si ces conditions sont vérifiées pour tous les littéraux concernés, alors
$C^\star=C_{\mathrm{true}}$ (égalité structurelle : $C^\star$ et
$C_{\mathrm{true}}$ désignent la même fonction booléenne, encodée par le
même ensemble de littéraux).
\end{theorem}

\begin{proof}
Pour $j\in J_+$ : $Y{=}1\Rightarrow X_j{=}1$, donc $C_j=P(X_j{=}0,Y{=}1)=0$,
et par \eqref{eq:pv} : $\Delta_j=S\,A_j-S\,B_j-a\cdot0=S(A_j-B_j)$, positif
ssi $A_j>B_j$. Par le Lemme~\ref{lem:excl}, $\bar\Delta_j=-a\,P(Y{=}1)-\Delta_j$ ;
comme $a\ge0$ et $P(Y{=}1)\ge0$, $\Delta_j>0\Rightarrow\bar\Delta_j\le-\Delta_j<0$.
Symétriquement pour $j\in J_-$ : $Y{=}1\Rightarrow X_j{=}0$ donne
$D_j=P(X_j{=}1,Y{=}1)=0$, d'où par \eqref{eq:pvbar} :
$\bar\Delta_j=S\,B_j-S\,A_j-a\cdot0=S(B_j-A_j)$, positif ssi $B_j>A_j$, et
alors $\Delta_j=-a\,P(Y{=}1)-\bar\Delta_j<0$.

Pour $k$ non pertinent, \eqref{eq:pv}--\eqref{eq:pvbar} donnent
$\Delta_k=S(A_k-B_k)-a\,C_k$ et $\bar\Delta_k=S(B_k-A_k)-a\,D_k$ (sans
simplification possible ici, $C_k,D_k$ quelconques). $\Delta_k<0\iff
A_k-B_k<\tfrac aS C_k$ ; $\bar\Delta_k<0\iff A_k-B_k>-\tfrac aS D_k$. Les deux
réunies donnent l'encadrement annoncé.
\end{proof}

\begin{remark}
Le Théorème~\ref{thm:identif} ne suppose ni indépendance ni équilibre des
features : $A_j,B_j,C_j,D_j$ sont des probabilités jointes quelconques et les
conditions ($A_j>B_j$, etc.) portent sur n'importe quelle loi jointe, pas sur
un modèle génératif particulier.
\end{remark}

\begin{theorem}[Robustesse à un bruit de label dépendant de la distribution]\label{thm:identif-bruit}
Reprenons les notations et hypothèses du Théorème~\ref{thm:identif}, et
notons $\Delta_j^\star,\bar\Delta_j^\star,\Delta_k^\star,\bar\Delta_k^\star$
les dérives calculées avec le label sans bruit $Y^\star$. Définissons
\[
  \mu = \min\Bigl(
    \{\Delta_j^\star : j\in J_+\} \cup \{-\bar\Delta_j^\star : j\in J_+\} \cup
    \{-\Delta_j^\star : j\in J_-\} \cup \{\bar\Delta_j^\star : j\in J_-\} \cup{}
  \]
  \[
    \{-\Delta_k^\star : k\notin J_+\cup J_-\} \cup \{-\bar\Delta_k^\star : k\notin J_+\cup J_-\}
  \Bigr),
\]
c'est-à-dire le plus petit des membres de gauche de toutes les inégalités
strictes exigées par le Théorème~\ref{thm:identif} (une quantité par
littéral et par signe requis) ; supposons $\mu>0$, i.e.\ ces inégalités
sont toutes satisfaites sans bruit. Soit maintenant $Y$ un label tel que
$\eta:=P(Y\ne Y^\star)>0$, \emph{sans aucune autre hypothèse} sur la loi du
bruit $Y\ne Y^\star$ (ni symétrie, ni indépendance vis-à-vis de $X$). Si
\[
  \boxed{\eta < \frac{\mu}{S+a},}
\]
alors toutes les inégalités du Théorème~\ref{thm:identif} restent
satisfaites (avec $Y$ à la place de $Y^\star$) et $C^\star=C_{\mathrm{true}}$.
Combiné au Théorème~\ref{thm:multi}, ceci donne directement
$\lim_{N\to\infty}\limsup_{t\to\infty}\mathbb P(C_t\not\equiv C_{\mathrm{true}})=0$ :
la règle récupère, avec probabilité tendant vers $1$, la conjonction
génératrice exacte -- pas seulement une structure définie implicitement par
$P$ -- tant que le bruit reste sous ce seuil.
\end{theorem}

\begin{proof}
\textbf{Étape 1 : chaque dérive est une espérance.}
Par \eqref{eq:pv}, $\Delta_j = S\,P(X_j{=}0,Y{=}0) - S\,P(X_j{=}1,Y{=}0) -
a\,P(X_j{=}0,Y{=}1)$. Ceci s'écrit $\Delta_j=\mathbb E[c(X_j,Y)]$ où
$c:\{0,1\}^2\to\mathbb R$ est définie par $c(0,0)=S$, $c(1,0)=-S$,
$c(0,1)=-a$, $c(1,1)=0$. De même, par \eqref{eq:pvbar},
$\bar\Delta_j=\mathbb E[\bar c(X_j,Y)]$ avec $\bar c(1,0)=S$,
$\bar c(0,0)=-S$, $\bar c(1,1)=-a$, $\bar c(0,1)=0$ (les rôles de $Y{=}0$ et
$Y{=}1$ sont échangés par rapport à $c$). Notons $\Delta_j^\star=\mathbb
E[c(X_j,Y^\star)]$ la même quantité calculée avec $Y^\star$ au lieu de $Y$.

\textbf{Étape 2 : un couplage naturel.}
Puisque $X$ suit la même loi dans les deux cas (seul le label diffère),
couplons $Y$ et $Y^\star$ sur le même espace de probabilité que $X$, de
sorte que $\{Y\ne Y^\star\}$ soit exactement l'événement de bruit, de
probabilité $\eta$. Un tel couplage existe toujours (c'est la définition
même de « $Y$ est un label bruité de $Y^\star$ avec taux $\eta$ »). Alors
\[
  \Delta_j - \Delta_j^\star = \mathbb E\bigl[c(X_j,Y) - c(X_j,Y^\star)\bigr]
  = \mathbb E\Bigl[\bigl(c(X_j,Y)-c(X_j,Y^\star)\bigr)\mathbf 1[Y\ne Y^\star]\Bigr],
\]
la dernière égalité car le terme entre crochets est nul dès que $Y=Y^\star$.

\textbf{Étape 3 : borne ponctuelle du saut de $c$, à $X_j$ fixé.}
Sur l'événement $Y\ne Y^\star$, les deux valeurs $Y,Y^\star$ sont
distinctes dans $\{0,1\}$, donc $\{Y,Y^\star\}=\{0,1\}$ dans un ordre ou
l'autre. Fixons la valeur de $X_j$ :
\begin{itemize}
  \item si $X_j=0$ : les deux valeurs possibles de $c(0,\cdot)$ sont
  $c(0,0)=S$ et $c(0,1)=-a$, donc $|c(0,Y)-c(0,Y^\star)|=|S-(-a)|=S+a$ ;
  \item si $X_j=1$ : les deux valeurs possibles sont $c(1,0)=-S$ et
  $c(1,1)=0$, donc $|c(1,Y)-c(1,Y^\star)|=|{-S}-0|=S$.
\end{itemize}
Dans les deux cas, $|c(X_j,Y)-c(X_j,Y^\star)|\le S+a$ (le maximum, $S+a$,
est atteint en $X_j=0$ ; puisque $a>0$, il domine $S$). Cette borne est
\emph{ponctuelle} : elle vaut pour toute réalisation de $X_j$, pas
seulement en moyenne.

\textbf{Étape 4 : passage à l'espérance.}
En combinant les étapes 2 et 3,
\[
  |\Delta_j-\Delta_j^\star| \le \mathbb E\Bigl[\bigl|c(X_j,Y)-c(X_j,Y^\star)\bigr|\,\mathbf 1[Y\ne Y^\star]\Bigr]
  \le (S+a)\,\mathbb E\bigl[\mathbf 1[Y\ne Y^\star]\bigr] = (S+a)\,\eta.
\]
Le même raisonnement appliqué à $\bar c$ (saut $S+a$ atteint en $X_j=1$
cette fois, par symétrie des rôles de $0,1$ sur $Y$) donne
$|\bar\Delta_j-\bar\Delta_j^\star|\le(S+a)\eta$. Le même argument, avec
exactement les mêmes fonctions $c,\bar c$, s'applique aux variables non
pertinentes $k$ et donne $|\Delta_k-\Delta_k^\star|\le(S+a)\eta$,
$|\bar\Delta_k-\bar\Delta_k^\star|\le(S+a)\eta$.

\textbf{Étape 5 : préservation des signes.}
Soit $\eta<\mu/(S+a)$, i.e.\ $(S+a)\eta<\mu$. Prenons par exemple
$j\in J_+$ : par définition de $\mu$, $\Delta_j^\star\ge\mu$. Par l'étape 4,
$\Delta_j\ge\Delta_j^\star-(S+a)\eta\ge\mu-(S+a)\eta>0$ puisque
$(S+a)\eta<\mu$. De même $-\bar\Delta_j^\star\ge\mu$ donne
$\bar\Delta_j\le\bar\Delta_j^\star+(S+a)\eta\le-\mu+(S+a)\eta<0$. Le même
calcul, terme à terme, s'applique à $j\in J_-$ et aux $k$ non pertinents.
Donc toutes les inégalités strictes du Théorème~\ref{thm:identif} restent
vraies avec $Y$ à la place de $Y^\star$, ce qui conclut.
\end{proof}

\begin{figure}[h]
\centering
\includegraphics[width=0.7\textwidth]{figures/graphe_eta_thm11.png}
\caption{Validation empirique du Théorème~\ref{thm:identif-bruit} : taux de
récupération exacte de $C_{\mathrm{true}}$ ($n{=}6$, $m{=}3$, $S{=}1$,
$a{=}0.3$) en fonction du bruit de label $\eta$, sur $200$ répétitions par
point. La transition empirique se produit nettement autour du seuil exact
$\eta^\star{=}2S/(4S+a\cdot2^m){=}0.3125$, calculé directement pour ce cas
i.i.d.\ Bernoulli(1/2) -- la borne générale du Théorème~\ref{thm:identif-bruit},
$\mu/(S{+}a)\approx0.014$, est confirmée très conservatrice pour ce cas
particulier.}
\label{fig:thm11}
\end{figure}

\section{Ensemble conditionné et boosté}

On intègre la règle à 4 cas dans un ensemble conditionné (conditions
aléatoires) et boosté (AdaBoost) pour illustrer son utilité pratique en
classification. \textbf{L'analyse de convergence complète de cet ensemble --
couverture des conditions sous la distribution adaptative $Q_m$, et
convergence du vote pondéré final -- n'est pas prouvée ici et reste un
travail futur.} Les résultats de cette section (couverture, boosting)
caractérisent des propriétés partielles utiles mais ne ferment pas cette
boucle.

\subsection{Conditions aléatoires}

Pour chaque clause $m$, on tire une condition $F_m\subseteq\{1,\ldots,n\}$,
$|F_m|=K$, et une configuration $z_m\in\{0,1\}^K$, uniformément
indépendamment. La clause n'est active que sur les exemples qui satisfont
sa condition :
\[
  g_m(x) = \mathbf 1[x_{F_m}=z_m].
\]
Les variables de $F_m$ sont retirées de l'apprentissage local : leurs deux
automates ($v_i$ et $\vbar_i$ pour $i\in F_m$) sont fixés dans l'état
Exclude et ne sont plus mis à jour -- la condition $x_{F_m}=z_m$ elle-même
n'est pas encodée par des littéraux inclus dans la clause, mais imposée
séparément par le mécanisme d'activation $g_m$. Seules les $n-K$ variables
restantes sont gouvernées par $2(n-K)$ automates entraînables par la règle
à quatre cas.

\subsection{Entraînement de la clause}

Pour chaque clause $m$ : (1) on sélectionne le sous-ensemble $S_m$ des
exemples d'entraînement satisfaisant $g_m$ ; (2) on équilibre le tirage
50/50 entre les deux classes lorsque les deux sont présentes dans $S_m$ ;
(3) au sein de chaque classe, le tirage est proportionnel aux poids AdaBoost
courants $w_i$ ; (4) chaque exemple ainsi tiré met à jour les $n-K$
automates libres par la règle à quatre cas.

\subsection{Agrégation pondérée}

Pour un classifieur binaire :
\[
  f(x) = \sum_{m=1}^{M} \alpha_m\, g_m(x)\,\bigl(2\,C_m(x)-1\bigr), \qquad
  \hat y = \mathbf 1[f(x)>0],
\]
avec
\[
  \alpha_m = \frac12\ln\frac{1-\mathrm{err}_m}{\mathrm{err}_m}.
\]
Chaque weak learner $h_m(x)=g_m(x)(2C_m(x)-1)\in\{-1,0,1\}$ s'abstient
(valeur $0$) lorsque sa condition n'est pas satisfaite par $x$ : il s'agit
donc d'un \emph{AdaBoost avec abstention}, où $\alpha_m$ et la mise à jour
des poids ne portent que sur la masse active $S_m$, et non d'AdaBoost
standard appliqué tel quel à un classifieur binaire toujours actif.

\subsection{Caractérisation par round de boosting}\label{sec:round-boosting}

Les Théorèmes~\ref{thm:stat}, \ref{thm:sel}, \ref{thm:multi},
\ref{thm:identif}, \ref{thm:identif-bruit} portent sur une clause unique
sous une loi $P(X,Y)$ \emph{fixe}. Le classifieur de la Section~5 entraîne
chaque clause $m$ sous une distribution adaptative $Q_m$ qui change à
chaque round de boosting -- ce résultat précise exactement ce qui, dans
cet écart, est couvert et ce qui ne l'est pas.

\begin{theorem}[Distribution d'entraînement effective d'un round]\label{thm:boosting}
Soit $S_m$ le sous-ensemble applicable de la clause $m$ (exemples
satisfaisant $g_m$) et $w_i$ les poids AdaBoost courants sur $S_m$.
Notons $W_m^+=\sum_{i\in S_m:y_i=1}w_i$, $W_m^-=\sum_{i\in S_m:y_i=0}w_i$.
L'Algorithme de la Section~5.6 ne tire \emph{pas} directement selon la
distribution pondérée $D_m(i)=w_i/\sum_{j\in S_m}w_j$ : il choisit d'abord
la classe positive ou négative avec probabilité $\tfrac12$ chacune si les
deux sont présentes dans $S_m$, sinon il échantillonne exclusivement dans
l'unique classe présente. La distribution effectivement échantillonnée
pour mettre à jour les automates de la clause $m$ est donc
\[
  Q_m(i) = \begin{cases}
    \tfrac12\,w_i/W_m^+, & i\in S_m,\ y_i{=}1,\ W_m^+,W_m^->0,\\
    \tfrac12\,w_i/W_m^-, & i\in S_m,\ y_i{=}0,\ W_m^+,W_m^->0,\\
    w_i/W_m^+, & i\in S_m,\ y_i{=}1,\ W_m^-=0,\\
    w_i/W_m^-, & i\in S_m,\ y_i{=}0,\ W_m^+=0,\\
    0, & i\notin S_m.
  \end{cases}
\]
Conditionnellement à l'historique du boosting et au tirage aléatoire de la
condition $(F_m,z_m)$ du round $m$, et sur l'évènement $S_m\ne\emptyset$,
les Théorèmes~\ref{thm:stat}, \ref{thm:sel}, \ref{thm:multi},
\ref{thm:identif}, \ref{thm:identif-bruit} s'appliquent \emph{tels quels} à
la clause $m$, chacun sous ses hypothèses respectives, avec $P(X,Y)$
remplacée par $Q_m(X,Y)$ -- \emph{et non} par $D_m$.
\end{theorem}

\begin{proof}
Les probabilités de transition \eqref{eq:pv}--\eqref{eq:pvbar} sont des
fonctions linéaires de la loi utilisée pour générer les tirages
effectivement présentés à l'automate ; elles ne supposent jamais une forme
particulière de cette loi, donc substituer $Q_m$ à $P$ ne change aucune
étape des preuves précédentes. $Q_m$, et non $D_m$, est la loi qui génère
les $(x,y)$ réellement vus par la mise à jour de la clause $m$ (tirage
50/50 puis proportionnel à $w$, Section~5.2) ; $D_m$ n'intervient que dans
le calcul de l'erreur pondérée et de $\alpha_m$ (Section~5.3).
\end{proof}

\begin{remark}[Ce que ce théorème ne prouve pas]
Il caractérise chaque round \emph{isolément}, conditionnellement à $Q_m$ --
il ne prouve pas la convergence du vote pondéré final
$f_c(x)=\sum_m\alpha_{c,m}h_{c,m}(x)$ vers un classifieur de Bayes global
(ce qu'AdaBoost classique établit sous des hypothèses différentes, non
reprises ici), ni que la structure identifiée round après round reste
stable ou cohérente d'un round à l'autre. Nous ne revendiquons pas ce
résultat plus fort : c'est précisément l'écart entre théorie de la règle
locale et convergence de l'ensemble complet annoncé en introduction et en
conclusion.
\end{remark}

\subsection{Extension multiclasse}

Pour chaque classe $c$, on pose $Y^{(c)}=\mathbf 1[Y{=}c]$ et on entraîne un
ensemble indépendant de clauses conditionnées et boostées pour ce problème
binaire :
\[
  f_c(x) = \sum_{m=1}^{M} \alpha_{c,m}\, g_{c,m}(x)\,\bigl(2\,C_{c,m}(x)-1\bigr),
  \qquad \hat y = \arg\max_c f_c(x).
\]

\subsection{Algorithme complet}

\begin{quote}\ttfamily\small
\textbf{Entrée} : dataset $\{(x_i,y_i)\}$, $n$ features, $N$ états, $M$
clauses, $K$ features gelées/clause, $total$ tirages/clause, $S,a$.\\[4pt]
$w_i \gets 1/|\text{train}|$ pour tout $i$\\[2pt]
\textbf{pour} $m = 1$ à $M$ :\\
\hspace*{1.5em}tirer $F\subseteq\{1,\ldots,n\}$, $|F|=K$, uniformément (shuffle)\\
\hspace*{1.5em}tirer $v\in\{0,1\}^K$ uniformément (pile-ou-face indépendant/feature)\\
\hspace*{1.5em}créer une clause $C_m$ conditionnée sur $(F,v)$ ; geler ces $K$ littéraux\\
\hspace*{1.5em}$S_m \gets \{i : (x_i)_F = v\}$\\
\hspace*{1.5em}\textbf{si} $S_m=\emptyset$ : $\alpha_m \gets 0$ ; \textbf{continuer}\\
\hspace*{1.5em}\textbf{pour} $t = 1$ à $total$ :\\
\hspace*{3em}tirer $i\in S_m$ selon $Q_m$ (50/50 classe, puis $\propto w_i$ dans la classe)\\
\hspace*{3em}mettre à jour $C_m$ avec $(x_i,\mathbf 1[y_i{=}c])$ (règle à 4 cas)\\
\hspace*{1.5em}\textbf{si} $C_m$ vide (aucun littéral inclus) : $\alpha_m \gets 0$ ; \textbf{continuer}
\hspace*{1em}(pas de mise à jour des poids)\\
\hspace*{1.5em}$\mathrm{err}_m \gets \sum_{i\in S_m} w_i\mathbf 1[C_m(x_i)\ne\mathbf 1[y_i{=}c]] \big/ \sum_{i\in S_m} w_i$ \hspace*{0.5em}(sous $D_m$)\\
\hspace*{1.5em}$\mathrm{err}_m \gets \mathrm{clip}(\mathrm{err}_m,\epsilon,1-\epsilon)$\\
\hspace*{1.5em}$\alpha_m \gets \tfrac12\ln\bigl((1-\mathrm{err}_m)/\mathrm{err}_m\bigr)$\\
\hspace*{1.5em}\textbf{pour} $i \in S_m$ : $w_i \gets w_i \cdot \exp(\pm\alpha_m)$ (signe $+$ si erreur, $-$ sinon)\\
\hspace*{1.5em}renormaliser $w$ sur tout le dataset\\[4pt]
\textbf{retourner} $(C_1,\ldots,C_M)$, $(\alpha_1,\ldots,\alpha_M)$
\end{quote}

\textbf{Prédiction.} Pour $x$ et chaque classe $c$ :
\[
  \mathrm{score}_c(x) = \sum_{m=1}^M \alpha_{c,m}\cdot\bigl(2\,C_{c,m}(x)-1\bigr)\cdot\mathbf 1\bigl[(x)_{F_{c,m}}=v_{c,m}\bigr]\cdot\mathbf 1[C_{c,m}\text{ non vide}],
\]
puis $\hat y = \arg\max_c \mathrm{score}_c(x)$ (clause ignorée si sa
condition n'est pas satisfaite par $x$ ou si elle est vide).

\subsection{Complexité}

Pour une classe donnée : la construction des sous-ensembles $S_m$ (tester,
pour chacun des $N_{\mathrm{train}}$ exemples et chacune des $M$ clauses, si
sa condition à $K$ features est satisfaite) coûte $O(M\,N_{\mathrm{train}}\,K)$ ;
l'entraînement proprement dit coûte ensuite $O(M\cdot total\cdot(n-K))$
(chacune des $M$ clauses reçoit $total$ tirages, chacun mettant à jour les
$n-K$ automates libres -- les $K$ automates gelés ne sont jamais touchés).
La mémoire est en $O(Mn)$ (deux automates par feature et par clause). La
prédiction sur un exemple coûte $O(Mn)$ dans le pire cas, mais en pratique
beaucoup moins : tester la condition de chacune des $M$ clauses coûte
$O(MK)$, et une condition à $K$ features tirée uniformément n'est
satisfaite que par une fraction $\approx2^{-K}$ des exemples en espérance,
de sorte que le nombre \emph{attendu} de clauses dont on évalue ensuite la
conjonction (coût $O(n-K)$ chacune) est $\approx M\cdot2^{-K}$. Le coût
moyen de prédiction est donc
\[
  O\bigl(MK + M\cdot2^{-K}\cdot(n-K)\bigr).
\]

\section{Validation empirique}

Comparaison de notre modèle à la Tsetlin Machine officielle
(\texttt{pyTsetlinMachine}, configurations publiées) et à cinq baselines
classiques (\texttt{scikit-learn}, hyperparamètres par défaut). Tous les
modèles sont évalués sur les mêmes fichiers de données et, quand un
remélange est nécessaire, la même seed par run ($20000+\text{run}$).
Écart-type rapporté : $\sqrt{\frac1R\sum_r(\mathrm{acc}_r-\bar{\mathrm{acc}})^2}$
(écart-type empirique de population, $1/R$, pas l'estimateur non biaisé
$1/(R{-}1)$). \textbf{Budgets non égalisés} entre notre modèle et la TM
(nombre de clauses différent par dataset, cf.\ ci-dessous, et même à nombre
de clauses égal -- cf.\ MNIST -- le nombre de mises à jour par clause
($total$) diffère du nombre d'epochs de la TM, donc le budget de calcul
proprement dit reste distinct) : ces comparaisons ne sont donc pas
équitables en termes de calcul. PCL (AAAI-25, related work le plus proche)
n'est pas comparé systématiquement sur XOR/Digits/MNIST : PCL s'effondre
sur XOR bruité avec ses hyperparamètres publiés (toutes les clauses restent
vides), et son propre article ne l'évalue lui-même que sur la récupération
de conjonctions synthétiques et sur Iris binaire -- pas sur des jeux
multiclasses larges. Nous le comparons donc directement là où la
comparaison est protocolairement fondée (Section~6.1 ci-dessous), en
réutilisant le code officiel de PCL.

\subsection{Comparaison directe à PCL}

En reprenant le protocole exact de l'article PCL (code officiel,
\texttt{dpcl-classifier/Main\_code\_AAAI25}), deux comparaisons directes :

\textbf{Récupération d'une conjonction cible.} Protocole et paramètres par
défaut du code officiel de PCL, inchangés : cible $C_{\mathrm{true}}$
tirée aléatoirement (chaque littéral indépendamment positif, négatif ou
absent, un littéral au moins forcé présent), table de vérité complète
($2^n$ exemples, ordre fixe, sans remélange), $100$ epochs (soit $100
\cdot 2^n$ mises à jour), $100$ runs, une seule clause, $N{=}10\,000$
états pour les deux méthodes, automates initialisés adjacents à la
frontière Include/Exclude ($N$ ou $N{+}1$). PCL : $p{=}0.75$ (défaut de
leur code). Notre règle : $S{=}1.0$, $a{=}0.3$. Succès = littéraux appris
exactement égaux aux littéraux cibles.
\begin{center}
\begin{tabular}{lccccc}
\toprule
$n$ & 4 & 6 & 8 & 10 & 12 \\
\midrule
PCL & 99/100 & 80/100 & 57/100 & 47/100 & 37/100 \\
Notre règle & \textbf{100/100} & \textbf{100/100} & \textbf{100/100} & \textbf{100/100} & \textbf{100/100} \\
\bottomrule
\end{tabular}
\end{center}
Ce protocole balaie la table de vérité dans un ordre fixe répété sur $100$
epochs -- une séquence déterministe périodique, pas un tirage i.i.d.\ de
$(X_t,Y_t)$ au sens du Théorème~\ref{thm:identif} ; sans bruit, ce balayage
complet représente néanmoins exactement la loi $P(X,Y)$ dégénérée voulue
(chaque combinaison de $X$ avec sa vraie étiquette, sans distorsion), donc
le résultat reste une validation directe et fidèle du théorème sur ce cas
limite. Notre règle récupère la cible exactement sur les $100$ cibles
testées à chaque $n$ (aucun échec observé, ce qui ne garantit pas une
probabilité de succès exactement égale à $1$), là où PCL se dégrade avec
$n$.

\textbf{Iris binaire.} Classe versicolor contre les deux autres, 16
features binaires, split 70/30 aléatoire, $p_i\sim U(0.6,0.8)$ par clause
pour PCL (16 clauses, 100 epochs) -- protocole et données de l'article PCL.
\begin{center}
\begin{tabular}{lc}
\toprule
Méthode & Précision \\
\midrule
PCL (code officiel, $92.4\%$ publié) & $92.22\%\pm3.80\%$ \\
Notre méthode complète ($M{=}300$, $K{=}2$, $total{=}2000$) & $93.60\%\pm3.63\%$ \\
\bottomrule
\end{tabular}
\end{center}
Notre reproduction de PCL ($92.22\%$) est cohérente avec leur chiffre
publié ($92.4\%$), validant le protocole. Les deux intervalles ($\pm1$
écart-type) se recouvrent largement et $M$ diffère (300 clauses contre 16) :
notre méthode complète est \emph{compétitive} avec PCL sur son propre
protocole Iris, sans qu'on puisse affirmer une supériorité statistiquement
significative sur cette seule expérience. Sur données réelles bruitées, la
règle seule (sans conditions ni boosting) dégénère en revanche nettement --
privée de tout mécanisme de diversité entre clauses, elle converge vers une
structure unique répétée $M$ fois ; c'est l'ensemble conditionné et boosté
complet (Section~5), et non la règle isolée, qui doit être comparé à PCL
sur ce type de données.

\subsection{Ablation : conditions et boosting}\label{sec:ablation}

Sur XOR bruité (configuration de la Section~6.3 : $M{=}200$, $K{=}1$,
$S{=}1.0$, $a{=}0.35$, $N{=}50$, $total{=}5000$), on active/désactive
indépendamment les conditions aléatoires et le boosting AdaBoost, à
hyperparamètres fixes ($10$ runs) :
\begin{center}
\begin{tabular}{lc}
\toprule
Variante & Précision \\
\midrule
A. Règle seule (sans conditions, sans boosting) & $48.86\%\pm0.00\%$ \\
B. + conditions seules (sans boosting) & $\mathbf{100.00\%\pm0.00\%}$ \\
C. + boosting seul (sans conditions) & $48.86\%\pm0.00\%$ \\
D. Modèle complet (conditions + boosting) & $\mathbf{100.00\%\pm0.00\%}$ \\
\bottomrule
\end{tabular}
\end{center}
Les variantes A et C sont numériquement identiques : sans conditions,
$S{=}1.0$ rend la mise à jour déterministe sur la branche dominante (les
exemples négatifs, majoritaires), si bien que toutes les clauses convergent
vers la même structure quels que soient les poids AdaBoost -- le
rééquilibrage des poids ne peut pas créer de diversité entre des clauses
qui ignorent ces poids. Sur ce problème, \textbf{l'expressivité provient
du conditionnement aléatoire, pas du boosting} : geler une variable rend
la seconde localement informative sur chaque sous-ensemble conditionné (un
mécanisme cohérent avec le Théorème~\ref{thm:sel}, mais dont la
contribution exacte à l'expressivité de l'ensemble n'est pas analysée
formellement, cf.\ Théorème~\ref{thm:boosting}). Le boosting n'ajoute rien
ici au-delà des conditions ($B=D$) ; son rôle -- pondérer et filtrer les
clauses faibles à plus grande échelle -- est plausiblement plus utile sur
des problèmes multiclasses (Digits, MNIST) que sur XOR, mais cela n'est pas
testé ici.

\subsection{XOR}
Fichiers fixes (\texttt{NoisyXORTrainingData.txt}/\texttt{TestData.txt},
5000/5000), aucun remélange. Notre modèle : $M{=}200$, $K{=}1$, $S{=}1.0$,
$a{=}0.35$, $N{=}50$, $total{=}5000$. TM : 20 clauses, $T{=}15$, $s{=}3.9$,
200 epochs. 20 runs partout.
\begin{center}
\begin{tabular}{lc}
\toprule
Modèle & Précision \\
\midrule
Notre modèle & $\mathbf{100\%\pm0\%}$ \\
Tsetlin Machine (officielle) & $97.40\%\pm2.23\%$ \\
SVM & $98.12\%\pm0\%$ \\
MLP & $76.72\%\pm1.97\%$ \\
Random Forest & $66.52\%\pm0.30\%$ \\
Régression logistique & $48.94\%\pm0\%$ \\
Naive Bayes & $49.12\%\pm0\%$ \\
\bottomrule
\end{tabular}
\end{center}

\subsection{Iris}
Fichier unique (\texttt{BinaryIrisData.txt}, 149 exemples), remélangé 80/20 à
chaque run. Notre modèle : $M{=}1500$, $K{=}2$, $S{=}1.0$, $a{=}0.3$,
$N{=}100$, $total{=}20000$. TM : 300 clauses, $T{=}10$, $s{=}3.0$, 500 epochs.
20 runs partout.
\begin{center}
\begin{tabular}{lc}
\toprule
Modèle & Précision \\
\midrule
Notre modèle & $\mathbf{95.5\%\pm4.25\%}$ \\
Random Forest & $95.33\%\pm3.86\%$ \\
MLP & $94.50\%\pm3.38\%$ \\
Tsetlin Machine (officielle) & $93.50\%\pm4.28\%$ \\
SVM & $92.67\%\pm4.03\%$ \\
Régression logistique & $92.50\%\pm4.70\%$ \\
Naive Bayes & $90.67\%\pm4.29\%$ \\
\bottomrule
\end{tabular}
\end{center}
Notre modèle et Random Forest sont dans la marge de bruit l'un de l'autre.

\subsection{Digits}
Fichiers fixes (\texttt{BinaryDigitsTrainingData192.txt}/\texttt{TestData192.txt},
1437/360), aucun remélange. Notre modèle : $M{=}2000$, $K{=}2$, $S{=}0.05$,
$a{=}1.0$, $N{=}100$, $total{=}20000$. TM : 1000 clauses, $T{=}10$, $s{=}3.0$,
300 epochs. 10 runs partout.
\begin{center}
\begin{tabular}{lc}
\toprule
Modèle & Précision \\
\midrule
Notre modèle & $\mathbf{97.19\%\pm0.40\%}$ \\
SVM & $97.50\%\pm0\%$ \\
Random Forest & $97.31\%\pm0.31\%$ \\
MLP & $97.17\%\pm0.41\%$ \\
Régression logistique & $96.11\%\pm0\%$ \\
Tsetlin Machine (officielle) & $93.69\%\pm0.76\%$ \\
Naive Bayes & $91.39\%\pm0\%$ \\
\bottomrule
\end{tabular}
\end{center}

\subsection{MNIST}
Fichiers fixes (\texttt{MNISTTraining.txt}/\texttt{Test.txt}, 60000/10000),
aucun remélange. Notre modèle : $M{=}16000$, $K{=}2$, $S{=}1.0$, $a{=}0.3$,
$N{=}100$, $total{=}20000$. TM : 4000 clauses/classe, $T{=}50$, $s{=}10.0$,
20 epochs.
\begin{center}
\begin{tabular}{lc}
\toprule
Modèle & Précision \\
\midrule
SVM & $\mathbf{97.71\%\pm0\%}$ \\
Random Forest & $97.07\%\pm0.06\%$ \\
MLP & $96.90\%\pm1.03\%$ \\
Tsetlin Machine (officielle, 4000 clauses/classe) & $96.62\%\pm0.07\%$ (3 runs) \\
Notre modèle ($M{=}16000$/classe) & $95.85\%$ (1 run) \\
Régression logistique & $91.89\%\pm0\%$ \\
Notre modèle ($M{=}4000$/classe, nombre de clauses égal à la TM) & $93.56\%$ (1 run) \\
Naive Bayes & $84.23\%\pm0\%$ \\
\bottomrule
\end{tabular}
\end{center}

\subsection{Effet de $M$ sur la précision}

Sur XOR (config.\ Section~6.1 fixée, $M$ variable, 20 runs/point) :
\begin{center}
\begin{tabular}{lccccccc}
\toprule
$M$ & 5 & 10 & 20 & 50 & 100 & 200 & 400 \\
\midrule
Précision & 64.9\% & 75.1\% & 84.7\% & 96.8\% & 99.3\% & 100\% & 100\% \\
\bottomrule
\end{tabular}
\end{center}
Rendements décroissants nets, saturés dès $M{=}200$ (configuration
retenue en Section~6.1). %TODO : ajouter Digits/MNIST des que le sweep
% Onyxia (M=100..2000 pour Digits, M=500..2000 pour MNIST) est termine ;
% points M=4000/8000/16000 deja disponibles ci-dessus pour MNIST.

\section{Limites et conclusion}

\textbf{Ce qui est prouvé.} La règle locale à quatre cas admet une loi
stationnaire exacte sous toute loi $P(X,Y)$ (Théorème~\ref{thm:stat}), ne
peut pas favoriser stationnairement l'inclusion simultanée d'un littéral et
de sa négation (Lemme~\ref{lem:excl} -- à $N$ fini, l'état instantané peut
néanmoins se trouver transitoirement du mauvais côté, avec probabilité
$\le\rho^{-N}\to0$), présente exactement l'un de trois régimes de
préférence stationnaire déterminés par le signe de $\Delta,\bar\Delta$
(Théorème~\ref{thm:sel}), et
concentre, pour une clause entière à $n-K$ variables libres, autour de la
structure stationnaire préférée $C^\star$ à un taux $2(n-K)\gamma_{\min}^{-N}$
(Théorème~\ref{thm:multi}). Sous des conditions explicites d'identifiabilité
(dépendance et déséquilibre arbitraires des features, Théorème~\ref{thm:identif}),
cette concentration porte sur la véritable conjonction génératrice
$C_{\mathrm{true}}$ et non plus seulement sur une structure définie
implicitement par $P$, et reste valide sous un bruit de labels arbitraire
de probabilité $\eta<\mu/(S+a)$ (Théorème~\ref{thm:identif-bruit}). Ces résultats sont chacun validés
empiriquement par simulation directe de la dynamique (Figures~\ref{fig:thm6}
et~\ref{fig:thm7}) et, pour l'identification, par comparaison directe à
PCL sur son propre protocole (Section~6.1).

\textbf{Ce qui ne l'est pas.} Le Théorème~\ref{thm:boosting} montre que
chaque round de boosting, pris isolément, est couvert exactement par les
résultats ci-dessus (avec $P$ remplacée par la distribution effective
$Q_m$) : l'écart restant n'est donc pas que « la théorie ne dit rien sur
le système entraîné », mais que la \emph{composition} entre rounds --
convergence du vote pondéré final $\sum_m\alpha_m h_m$ vers un classifieur
global -- n'est pas analysée. Les résultats sur MNIST à $M$ élevé ($8000$,
$16000$) reposent sur un seul run chacun, sans écart-type ; les
comparaisons à la TM ne sont pas systématiquement à budget de calcul égal ;
PCL n'est pas inclus comme baseline complète en dehors d'Iris et de XOR ;
l'ablation de la Section~6.2 isole l'effet des conditions et du boosting
sur XOR, mais une étude systématique de $S,a,N,K,M,total$ pris séparément
n'a pas été menée à ce stade.

\textbf{Conclusion.} En découplant entièrement la mise à jour de chaque
automate de la sortie de la clause et du vote de l'ensemble, la règle à
quatre cas admet une analyse de convergence exacte, sous une loi $P(X,Y)$
arbitraire, que ni les machines de Tsetlin ni PCL ne permettent sous cette
généralité. Intégrée dans un ensemble conditionné et boosté simple, elle
atteint $100\%\pm0\%$ sur XOR bruité et des performances compétitives sur
Iris, Digits et MNIST, tout en laissant ouverte l'analyse de convergence de
l'ensemble complet.

\end{document}
